% !TeX root = ../../../Book1.tex
\section{\texorpdfstring{\(\pi\)}{π}-系统、\texorpdfstring{\(\lambda\)}{λ}-系统与单调类}
\label{b1:sec:0103}

若两个测度或两个积分公式都定义在一个巨大的 \(\sigma\)-代数上，逐个可测集验证它们相同通常没有可操作性。本节的基本策略是：先选一组生成元，在其上验证目标性质，再利用目标性质本身的闭性把结论推广出去。

\subsection{\texorpdfstring{\(\pi\)}{π}-系统}
\label{b1:subsec:010301}

\begin{definition}
\label{b1:def:pi-system}
集合族 \(\mathcal P\subseteq\mathcal P(X)\) 若对任意 \(A,B\in\mathcal P\) 都有 \(A\cap B\in\mathcal P\)，则称为 \term[pi-xitong]{\(\pi\)-系统}[pi-system]。
\end{definition}

\(\pi\)-系统只要求有限交封闭，不要求补集或并封闭，因而比 \(\sigma\)-代数容易识别。区间、半直线和矩形族都是典型来源。

\begin{remark}
\label{b1:rem:pi-systems}
\textbf{例。} 在 \(\mathbb R\) 上，半直线族 \(\{(-\infty,t]:t\in\mathbb R\}\) 是 \(\pi\)-系统；在 \(X\times Y\) 上，所有可测矩形 \(A\times B\) 也构成 \(\pi\)-系统，因为
\[
(A_1\times B_1)\cap(A_2\times B_2)
=(A_1\cap A_2)\times(B_1\cap B_2).
\]
\end{remark}

\subsection{Dynkin \texorpdfstring{\(\lambda\)}{λ}-系统}
\label{b1:subsec:010302}

\begin{definition}
\label{b1:def:lambda-system}
\term[dynkin-lambda-xitong]{Dynkin \(\lambda\)-系统}[Dynkin lambda-system] 是集合族 \(\mathcal D\subseteq\mathcal P(X)\)，满足：
\begin{enumerate}
  \item \(X\in\mathcal D\)；
  \item 若 \(A,B\in\mathcal D\) 且 \(A\subseteq B\)，则 \(B\setminus A\in\mathcal D\)；
  \item 若 \(A_1,A_2,\ldots\in\mathcal D\) 两两不交，则 \(\bigcup_{n=1}^{\infty}A_n\in\mathcal D\)。
\end{enumerate}
\end{definition}

第二条取 \(B=X\) 立即给出补集封闭性。每个 \(\sigma\)-代数都是 \(\lambda\)-系统；反之，\(\lambda\)-系统未必对交封闭。对集合族 \(\mathcal E\)，记 \(\lambda(\mathcal E)\) 为包含 \(\mathcal E\) 的最小 \(\lambda\)-系统，它可定义为所有包含 \(\mathcal E\) 的 \(\lambda\)-系统之交。

\subsection{\texorpdfstring{\(\pi\)}{π}-\texorpdfstring{\(\lambda\)}{λ} 定理}
\label{b1:subsec:010303}

\begin{theorem}
\label{b1:thm:pi-lambda}
若 \(\mathcal P\) 是 \(X\) 上的 \(\pi\)-系统，则
\[
\lambda(\mathcal P)=\sigma(\mathcal P).
\]
因而，任何包含 \(\mathcal P\) 的 \(\lambda\)-系统都包含 \(\sigma(\mathcal P)\)。
\end{theorem}
\begin{proof}
令 \(\mathcal D=\lambda(\mathcal P)\)。先定义
\[
\mathcal D_0
=\{A\in\mathcal D:A\cap P\in\mathcal D
  \text{ 对每个 }P\in\mathcal P\text{ 成立}\}.
\]
利用交与差、两两不交并可以逐项验证 \(\mathcal D_0\) 是 \(\lambda\)-系统。又因 \(\mathcal P\) 对交封闭，\(\mathcal P\subseteq\mathcal D_0\)。由 \(\mathcal D\) 的最小性，\(\mathcal D=\mathcal D_0\)。

现在固定 \(A\in\mathcal D\)，并令
\[
\mathcal D_A=\{B\in\mathcal D:A\cap B\in\mathcal D\}.
\]
同样地，\(\mathcal D_A\) 是 \(\lambda\)-系统。刚证明的 \(\mathcal D=\mathcal D_0\) 表明 \(\mathcal P\subseteq\mathcal D_A\)，故 \(\mathcal D_A=\mathcal D\)。这说明 \(\mathcal D\) 对有限交封闭。

\(\lambda\)-系统一旦对有限交封闭，便对有限并封闭；再把任意可数并写成两两不交部分的可数并，可知 \(\mathcal D\) 是 \(\sigma\)-代数。因此 \(\sigma(\mathcal P)\subseteq\mathcal D\)。反向包含成立是因为每个 \(\sigma\)-代数都是 \(\lambda\)-系统，故等式成立。
\end{proof}

\begin{corollary}
\label{b1:cor:measure-uniqueness-pi}
设 \(\mathcal P\) 是含 \(X\) 的 \(\pi\)-系统，\(\mu,\nu\) 是 \(\sigma(\mathcal P)\) 上的概率测度。若 \(\mu(P)=\nu(P)\) 对每个 \(P\in\mathcal P\) 成立，则 \(\mu=\nu\)。
\end{corollary}
\begin{proof}
令 \(\mathcal D=\{A\in\sigma(\mathcal P):\mu(A)=\nu(A)\}\)。由概率测度的可列可加性和 \(\mu(X)=\nu(X)=1\)，\(\mathcal D\) 是包含 \(\mathcal P\) 的 \(\lambda\)-系统。应用 \cref{b1:thm:pi-lambda} 即得 \(\mathcal D=\sigma(\mathcal P)\)。
\end{proof}

\subsection{单调类定理}
\label{b1:subsec:010304}

\begin{definition}
\label{b1:def:monotone-class}
集合族 \(\mathcal M\subseteq\mathcal P(X)\) 若满足：当 \(A_n\uparrow A\) 或 \(A_n\downarrow A\)，且每个 \(A_n\in\mathcal M\) 时，均有 \(A\in\mathcal M\)，则称 \(\mathcal M\) 为\term[dantiaolei]{单调类}[monotone class]。这里 \(A_n\uparrow A\) 表示 \(A_1\subseteq A_2\subseteq\cdots\) 且 \(\bigcup_nA_n=A\)；递减情形类似。
\end{definition}

\begin{theorem}
\label{b1:thm:monotone-class}
若 \(\mathcal A\) 是 \(X\) 上的集合代数，\(\mathcal M(\mathcal A)\) 是包含 \(\mathcal A\) 的最小单调类，则
\[
\mathcal M(\mathcal A)=\sigma(\mathcal A).
\]
\end{theorem}
\begin{proof}
记 \(\mathcal M=\mathcal M(\mathcal A)\)。先令
\[
\mathcal C=\{E\in\mathcal M:E^{\mathrm c}\in\mathcal M\}.
\]
对单调极限取补会把递增列变为递减列，反之亦然，因此 \(\mathcal C\) 是单调类。集合代数对补封闭，故 \(\mathcal A\subseteq\mathcal C\)，由最小性知 \(\mathcal C=\mathcal M\)。于是 \(\mathcal M\) 对补封闭。

固定 \(A\in\mathcal A\)，集合族
\[
\mathcal M_A=\{E\in\mathcal M:A\cap E\in\mathcal M\}
\]
是包含 \(\mathcal A\) 的单调类，故等于 \(\mathcal M\)。再固定任意 \(E\in\mathcal M\)，同理可得 \(\{F\in\mathcal M:E\cap F\in\mathcal M\}=\mathcal M\)。因此 \(\mathcal M\) 对有限交封闭，继而对有限并封闭。把可数并写成递增的有限并之极限，得到 \(\mathcal M\) 是 \(\sigma\)-代数。于是 \(\sigma(\mathcal A)\subseteq\mathcal M\)。

另一方面，\(\sigma(\mathcal A)\) 自身是包含 \(\mathcal A\) 的单调类，故 \(\mathcal M\subseteq\sigma(\mathcal A)\)。两边相等。
\end{proof}

\subsection{函数型单调类定理}
\label{b1:subsec:010305}

\begin{theorem}
\label{b1:thm:functional-monotone-class}
设 \(\mathcal A\) 是 \(X\) 上的集合代数，\(\mathcal H\) 是一族有界实值函数，满足：
\begin{enumerate}
  \item 常值函数属于 \(\mathcal H\)；
  \item \(\mathcal H\) 是实向量空间，且 \(h\in\mathcal H\) 时 \(\lvert h\rvert\in\mathcal H\)；
  \item 若 \(0\leq h_n\uparrow h\)，各 \(h_n\in\mathcal H\)，且 \(h\) 有界，则 \(h\in\mathcal H\)；
  \item 对每个 \(A\in\mathcal A\)，\(\mathbf 1_A\in\mathcal H\)。
\end{enumerate}
则每个有界的 \(\sigma(\mathcal A)\)-可测实值函数都属于 \(\mathcal H\)。
\end{theorem}
\begin{proof}
令 \(\mathcal D=\{E\subseteq X:\mathbf 1_E\in\mathcal H\}\)。由常值函数、向量空间性质和单调收敛封闭性，\(\mathcal D\) 是包含 \(\mathcal A\) 的单调类：递减列可先对 \(1-\mathbf 1_E\) 使用递增封闭性。由 \cref{b1:thm:monotone-class}，\(\mathcal D\) 包含 \(\sigma(\mathcal A)\)。

所以全部 \(\sigma(\mathcal A)\)-可测简单函数属于 \(\mathcal H\)。若 \(f\) 是有界非负可测函数，取逐点递增到 \(f\) 的有界简单函数列即可推出 \(f\in\mathcal H\)；一般有界实值函数写成正部减去负部即可。
\end{proof}

函数型定理的价值在于，它把集合上的推广原理直接转化为函数上的推广原理。后面验证积分恒等式时，常先对指标函数成立，再借本定理推广至所有有界可测函数。
